\chapter{LEAN\_Warmup: cut-transposes and the alternative theorem}

This chapter follows \texttt{LEAN\_Warmup.CutTranspose}, an outline of
Chatterjee--Ghosh--Gurjar--Raj, \emph{Characterizing and Testing Principal Minor Equivalence
of Matrices} (STOC 2025, full version arXiv:2410.01961), abbreviated CGGR. Its Theorem 1.1
is the second formalization target, and the alternative form of Theorem 4.1 is a corollary.
Throughout, $F$ is a field, $I$ is a finite index type and $X\subseteq I$. Results cited from
other papers are marked \emph{external}; the remaining unproved statements are marked
\emph{to do}. The algorithmic results of CGGR (Theorems 1.2 and 1.3), Lemma 2.4,
Observation 3.7, Claim 3.8 and the first part of Lemma 3.3 are not stated.

\section{Principal minor equivalence}

\begin{definition}[Principal minor equivalence]
\label{def:pme}
\lean{Matrix.PMEquiv, Matrix.pmEquiv_iff_mem_pmFiber, Matrix.PMEquiv.refl, Matrix.PMEquiv.symm, Matrix.PMEquiv.trans}
\leanok
\uses{def:fiber}
Matrices $A,B$ over a commutative ring are \emph{principal minor equivalent} (PME) if
$A_S=B_S$ for every finite $S$. Equivalently $B$ lies in the principal minor fiber of $A$.
The relation is reflexive, symmetric and transitive.
\end{definition}

\begin{lemma}[Transpose and diagonal equivalence]
\label{lem:pme-basic}
\lean{Matrix.pmEquiv_transpose, Matrix.PMEquiv.transpose, Matrix.DiagEquivalent.pmEquiv, Matrix.DiagSimilar.pmEquiv, Matrix.hasTrivialFiber_iff_forall_pmEquiv}
\leanok
\uses{def:pme,lem:diag-principal-minor}
A matrix is PME to its transpose, and PME is preserved by transposing both matrices.
Diagonally equivalent matrices are PME. The fiber of $A$ is trivial exactly when every
matrix PME to $A$ is diagonally equivalent to $A$.
\end{lemma}
\begin{proof}
\leanok
Transposition preserves determinants; the rest restates Lemma~\ref{lem:diag-principal-minor}.
\end{proof}

\begin{lemma}[Principal submatrices]
\label{lem:pme-submatrix}
\lean{Matrix.PMEquiv.principalSubmatrix}
\leanok
\uses{def:pme}
If $A$ and $B$ are PME, so are $A[T,T]$ and $B[T,T]$ for every finite $T$.
\end{lemma}
\noindent\emph{Status: to do. A principal minor of $A[T,T]$ at $S\subseteq T$ is the principal
minor of $A$ at the image of $S$.}

\begin{lemma}[PME and determinantal polynomials]
\label{lem:pme-detpoly}
\lean{Matrix.pmEquiv_iff_detpoly_eq}
\leanok
\uses{def:pme,thm:detpoly-expansion}
$A$ and $B$ are PME exactly when $f_A=f_B$.
\end{lemma}
\begin{proof}
\leanok
This is Theorem~\ref{thm:detpoly-expansion}.
\end{proof}

\begin{definition}[Nonzero off-diagonal entries]
\label{def:offdiag}
\lean{Matrix.OffDiagNeZero, Matrix.OffDiagNeZero.isIrred, Matrix.OffDiagNeZero.transpose}
\leanok
\uses{lem:reducible}
$A$ has \emph{nonzero off-diagonal entries} if $A_{ij}\ne0$ for all $i\ne j$. Such a matrix
is irreducible, and the condition is preserved by transpose. CGGR reduce Theorem 1.1 to
this case.
\end{definition}

\begin{lemma}[PME preserves irreducibility]
\label{lem:pme-irred}
\lean{Matrix.PMEquiv.isIrred}
\leanok
\uses{lem:pme-detpoly,lem:irred-poly}
For nonempty $I$, if $A$ and $B$ are PME and $A$ is irreducible, then $B$ is irreducible.
\end{lemma}
\begin{proof}
\leanok
By Lemma~\ref{lem:irred-poly}, irreducibility of the matrix is irreducibility of $f_A=f_B$.
\end{proof}

\section{Cuts and minimal cuts}

\begin{definition}[Minimal cuts]
\label{def:minimal-cut}
\lean{Matrix.IsMinimalCut, Matrix.HasCut.isMinimalCut_of_card_eq_two}
\leanok
\uses{def:cut}
A cut $X$ of $A$ is \emph{minimal} if no cut of $A$ is properly contained in $X$.
A cut of size two is minimal.
\end{definition}

\begin{lemma}[Existence of minimal cuts]
\label{lem:minimal-cut-exists}
\lean{Matrix.HasCut.exists_isMinimalCut}
\leanok
\uses{def:minimal-cut}
Every cut contains a minimal cut.
\end{lemma}
\noindent\emph{Status: to do. Well-founded induction on strict inclusion.}

\section{The cut-transpose operation}

The definitions and the basic API are in Section~\ref{sec:cuts} of Part~\ref{part:detpoly}.
This section adds what is only stated over a field.

\begin{lemma}[Uniqueness up to diagonal similarity (CGGR Remark 2.11)]
\label{lem:ct-unique}
\lean{Matrix.IsCutTranspose.diagSimilar}
\leanok
\uses{def:cut-transpose,lem:cut-ranks,def:diag}
If $A$ is irreducible with a cut $X$, any two cut-transposes of $A$ at $X$ are diagonally
similar.
\end{lemma}
\noindent\emph{Status: to do. The cross blocks have rank exactly one
(Lemma~\ref{lem:cut-ranks}), so two factorizations differ by
$(\lambda p,\lambda^{-1}q,\mu u,\mu^{-1}v)$; conjugate by the diagonal matrix that is
$\lambda\mu$ on $X$ and $1$ on $X^c$.}

\begin{lemma}[Cut-transpose preserves principal minors (CGGR Lemma 2.12)]
\label{lem:ct-pme}
\lean{Matrix.cutTranspose_eq_cutSwap, Matrix.IsCutTranspose.pmEquiv}
\leanok
\uses{lem:ct-principal-minor,def:pme,def:cut-swap}
If $B$ is a cut-transpose of $A$ at $X$, then $A$ and $B$ are PME. This is the matrix form of
Lemma 4.5 of the paper: the cut-transpose is the cut swap of Definition~\ref{def:cut-swap}
applied to $A$ instead of to $G_A$.
\end{lemma}
\begin{proof}
\leanok
Lemma~\ref{lem:ct-principal-minor}.
\end{proof}

\begin{lemma}[Compatibility with diagonal equivalence]
\label{lem:ct-de-compat}
\lean{Matrix.DiagEquivalent.exists_isCutTranspose}
\leanok
\uses{lem:ct-transpose,def:diag}
If $A'$ is diagonally equivalent to $A$ and $B'$ is a cut-transpose of $A'$ at $X$, then
$B'$ is diagonally equivalent to some cut-transpose of $A$ at $X$.
\end{lemma}
\noindent\emph{Status: to do. Use Lemma~\ref{lem:ct-transpose} for the transpose, and
$\operatorname{ct}_X(DAD^{-1})=\tilde D\operatorname{ct}_X(A)\tilde D^{-1}$ with
$\tilde D=\operatorname{diag}(D_X,D_{X^c}^{-1})$ for the induced factorization.}

\begin{lemma}[Diagonal equivalence preserves cuts]
\label{lem:de-cuts}
\lean{Matrix.DiagEquivalent.hasCut}
\leanok
\uses{def:cut,def:diag}
Diagonally equivalent matrices have the same cuts.
\end{lemma}
\noindent\emph{Status: to do. Rescale a cut factorization, or use that multiplication by an
invertible diagonal matrix preserves rank.}

\section{Cut-transpose equivalence}

\begin{definition}[Chains and cut-transpose equivalence (CGGR Definition 2.14)]
\label{def:ct-chain}
\lean{Matrix.CutTransposeEquivWrt, Matrix.CutTransposeEquiv, Matrix.cutTransposeEquiv_iff_exists_cutTransposeEquivWrt, Matrix.DiagEquivalent.cutTransposeEquiv}
\leanok
\uses{def:ct-chain-detpoly,def:diag}
With cut-transpose paths as in Definition~\ref{def:ct-chain-detpoly}, $A$ and $B$ are
\emph{cut-transpose equivalent with respect to} a list of cuts if a path from $A$ with these
cuts ends at a matrix diagonally equivalent to $B$, and \emph{cut-transpose equivalent} if
some path from $A$ ends at a matrix diagonally equivalent to $B$; equivalently, if the first
holds for some list. Diagonally equivalent matrices are cut-transpose equivalent through the
empty path.
\end{definition}

\begin{lemma}[Theorem 1.1, easy direction]
\label{lem:chain-pme}
\lean{Matrix.CutTransposePath.pmEquiv, Matrix.CutTransposeEquivWrt.pmEquiv, Matrix.CutTransposeEquiv.pmEquiv}
\leanok
\uses{def:ct-chain,lem:chain-principal-minor,lem:pme-basic}
Cut-transpose equivalent matrices are PME.
\end{lemma}
\begin{proof}
\leanok
Lemma~\ref{lem:chain-principal-minor}, then Lemma~\ref{lem:pme-basic}.
\end{proof}

\begin{lemma}[Chains inside a diagonal-equivalence class]
\label{lem:chain-de}
\lean{Matrix.CutTransposePath.diagEquivalent}
\leanok
\uses{def:ct-chain,lem:ct-de-compat,lem:de-cuts}
Suppose every cut-transpose of $A$ at every cut of $A$ is diagonally equivalent to $A$.
Then every path starting at a matrix diagonally equivalent to $A$ ends at a matrix
diagonally equivalent to $A$.
\end{lemma}
\begin{proof}
\leanok
Induct on the path. At each step transfer the cut and the cut-transpose to $A$ by
Lemmas~\ref{lem:de-cuts} and~\ref{lem:ct-de-compat}.
\end{proof}

\section{Passing to the adjugate}

These statements reduce Theorem 1.1 to matrices with nonzero off-diagonal entries.
CGGR's ``$A+D$ is nonsingular'' is the hypothesis that $\det(A+\operatorname{diag}(d))$ is
a unit.

\begin{lemma}[PME and adjugates (CGGR Lemma 2.1)]
\label{ext:cggr21}
\lean{Matrix.pmEquiv_adjugate_add_diagonal_iff}
\leanok
\uses{def:pme}
If $A+D$ and $B+D$ are nonsingular for a diagonal $D$, then
$\operatorname{adj}(A+D)$ and $\operatorname{adj}(B+D)$ are PME exactly when $A$ and $B$ are.
\end{lemma}
\noindent\emph{Status: external, from Hartfiel--Loewy (1984), Lemma 4.}

\begin{lemma}[Cuts of the adjugate (CGGR Lemma 2.6)]
\label{lem:cggr26}
\lean{Matrix.hasCut_adjugate_add_diagonal_iff}
\leanok
\uses{def:cut}
If $A+D$ is nonsingular, $A$ and $\operatorname{adj}(A+D)$ have the same cuts.
\end{lemma}
\noindent\emph{Status: to do. Lemma~\ref{lem:adj-cut} is the same statement before
specializing $x\mapsto d$.}

\begin{lemma}[Diagonal similarity and adjugates (CGGR Lemma 2.9)]
\label{ext:cggr29}
\lean{Matrix.diagSimilar_adjugate_add_diagonal_iff}
\leanok
\uses{def:diag}
If $A+D$ and $B+D$ are nonsingular, then $\operatorname{adj}(A+D)$ and
$\operatorname{adj}(B+D)$ are diagonally similar exactly when $A$ and $B$ are.
\end{lemma}
\noindent\emph{Status: external, from Hartfiel--Loewy (1984), Lemma 4.}

\begin{lemma}[Cut-transpose commutes with the adjugate (CGGR Lemma 2.13)]
\label{lem:cggr213}
\lean{Matrix.IsCutTranspose.adjugate}
\leanok
\uses{def:cut-transpose,lem:ct-unique}
If $A$ is irreducible with a cut $X$ and $B$ is a cut-transpose of $A$ at $X$, then $X$ is
a cut of $\operatorname{adj}(A)$ and $\operatorname{adj}(B)$ is diagonally similar to a
cut-transpose of $\operatorname{adj}(A)$ at $X$.
\end{lemma}
\noindent\emph{Status: to do.}

\begin{lemma}[Cut-transpose equivalence and adjugates (CGGR Lemma 2.15)]
\label{lem:cggr215}
\lean{Matrix.cutTransposeEquivWrt_adjugate_add_diagonal_iff}
\leanok
\uses{def:ct-chain,lem:cggr26,ext:cggr29,lem:cggr213}
For irreducible $A,B$ with $A+D$ and $B+D$ nonsingular, $\operatorname{adj}(A+D)$ and
$\operatorname{adj}(B+D)$ are cut-transpose equivalent with respect to a list exactly when
$A$ and $B$ are.
\end{lemma}
\noindent\emph{Status: to do.}

\begin{lemma}[A good diagonal shift exists (CGGR Claim 4.1)]
\label{lem:claim41}
\lean{Matrix.exists_diagonal_offDiagNeZero_adjugate}
\leanok
\uses{ext:nonzero,def:offdiag}
Over an infinite field, for irreducible $A,B$ there is a diagonal $D$ with $A+D$, $B+D$
nonsingular and all entries of both adjugates nonzero. CGGR require only $|F|>10n^5$.
\end{lemma}
\noindent\emph{Status: to do. Every entry of $G_A$ is a nonzero polynomial
(Lemma~\ref{ext:nonzero}), and a finite product of nonzero polynomials over an infinite
field has a non-root.}

\section{The characterization}

\begin{lemma}[Small matrices (CGGR Lemma 2.8 (1))]
\label{ext:hl-small}
\lean{Matrix.hasTrivialFiber_of_card_le_three}
\leanok
\uses{def:fiber,def:irred}
An irreducible matrix with $|I|\le3$ has a trivial fiber. Part (2) of the lemma is
Lemma~\ref{ext:no-cut}.
\end{lemma}
\noindent\emph{Status: external, Hartfiel--Loewy (1984), Theorem 3.}

\begin{lemma}[The case $|I|=4$ (CGGR Lemma 3.1)]
\label{lem:cggr31}
\lean{Matrix.diagEquivalent_or_cutTranspose_of_card_eq_four}
\leanok
\uses{def:offdiag,def:pme,def:cut-transpose}
Let $|I|=4$, let $A$ have nonzero off-diagonal entries and let $B$ be PME to $A$. Then
$A$ and $B$ are diagonally equivalent, or they have a common cut and every cut-transpose
of $A$ at a common cut is diagonally equivalent to $B$.
\end{lemma}
\noindent\emph{Status: to do.}

\begin{lemma}[Cuts of a cut-transpose (CGGR Lemma 3.2)]
\label{lem:cggr32}
\lean{Matrix.IsCutTranspose.hasCut_iff_of_nested, Matrix.IsCutTranspose.hasCut_symmDiff_iff}
\leanok
\uses{def:offdiag,def:cut-transpose}
Let $A$ have nonzero off-diagonal entries, $S$ a cut of $A$ and $C$ a cut-transpose of $A$
at $S$. If $T$ or $T^c$ is contained in $S$ or $S^c$, then $T$ is a cut of $C$ exactly when
it is a cut of $A$. Otherwise $T\mathbin{\triangle}S$ is a cut of $C$ exactly when $T$ is a
cut of $A$.
\end{lemma}
\noindent\emph{Status: to do.}

\begin{lemma}[Minimal cuts and one extra index (CGGR Lemma 3.3)]
\label{lem:cggr33}
\lean{Matrix.IsMinimalCut.not_hasCut_principalSubmatrix_insert}
\leanok
\uses{def:offdiag,def:minimal-cut}
If $A$ has nonzero off-diagonal entries, $S$ is a minimal cut with $|S|>2$ and $t\notin S$,
then $A[S\cup\{t\}]$ has no cut. Only this consequence of the lemma is stated.
\end{lemma}
\noindent\emph{Status: to do.}

\begin{lemma}[Lifting a cut (CGGR Lemma 3.4)]
\label{lem:cggr34}
\lean{Matrix.HasCut.of_principalSubmatrix_insert_compl}
\leanok
\uses{def:offdiag,def:cut}
If $A$ has nonzero off-diagonal entries, $S$ is a cut of $A$, $t\in S$ and $Y\subseteq S^c$
is a cut of $A[S^c\cup\{t\}]$, then $Y$ is a cut of $A$.
\end{lemma}
\noindent\emph{Status: to do.}

\begin{lemma}[Cuts of size two (CGGR Lemma 3.5)]
\label{lem:cggr35}
\lean{Matrix.PMEquiv.hasCut_or_of_card_eq_two}
\leanok
\uses{def:offdiag,def:pme,def:cut-transpose,lem:cggr31,lem:pme-submatrix}
Let $A,B$ be PME with nonzero off-diagonal entries and $S$ a cut of $A$ of size two.
Then $S$ is a cut of $B$, or $B$ has a cut $X$ such that $S$ is a cut of every
cut-transpose of $B$ at $X$. This is the form used in the proof of Theorem 1.1.
\end{lemma}
\noindent\emph{Status: to do.}

\begin{lemma}[Large minimal cuts (CGGR Lemma 3.6)]
\label{lem:cggr36}
\lean{Matrix.PMEquiv.hasCut_of_isMinimalCut}
\leanok
\uses{def:offdiag,def:pme,lem:cggr33,lem:pme-submatrix,ext:no-cut}
Let $A,B$ be PME with nonzero off-diagonal entries. A minimal cut of $A$ of size greater
than two is a cut of $B$.
\end{lemma}
\noindent\emph{Status: to do.}

\begin{lemma}[Gluing along a minimal cut (CGGR Lemma 3.9)]
\label{lem:cggr39}
\lean{Matrix.PMEquiv.diagSimilar_or_cutTranspose}
\leanok
\uses{def:offdiag,def:pme,def:minimal-cut,def:cut-transpose}
Let $A,B$ be PME with nonzero off-diagonal entries, $S$ a minimal cut of $A$ and a cut of
$B$, and $s\in S$. If $A[S^c\cup\{s\}]$ and $B[S^c\cup\{s\}]$ are diagonally similar, then
$A$ is diagonally similar to $B$, or every cut-transpose of $A$ at $S$ is.
\end{lemma}
\noindent\emph{Status: to do.}

\begin{theorem}[CGGR Theorem 1.1, hard direction]
\label{thm:cggr11-hard}
\lean{Matrix.PMEquiv.exists_cutTransposeEquivWrt}
\leanok
\uses{def:ct-chain,ext:cggr21,lem:cggr215,lem:claim41,ext:hl-small,ext:no-cut,lem:cggr31,lem:cggr32,lem:cggr34,lem:cggr35,lem:cggr36,lem:cggr39,lem:minimal-cut-exists,lem:pme-submatrix}
Irreducible PME matrices are cut-transpose equivalent with respect to a list of length at
most $2|I|$.
\end{theorem}
\noindent\emph{Status: to do; a formalization target. Reduce to nonzero off-diagonal
entries by passing to $\operatorname{adj}(A+D)$, possibly over a field extension, then
induct on $|I|$ through a minimal cut $S$ and the principal submatrix on $S^c\cup\{s\}$.
The base cases are Lemmas~\ref{ext:hl-small}, \ref{ext:no-cut} and~\ref{lem:cggr31}.}

\begin{theorem}[CGGR Theorem 1.1]
\label{thm:cggr11}
\lean{Matrix.pmEquiv_iff_cutTransposeEquivWrt, Matrix.pmEquiv_iff_cutTransposeEquiv}
\leanok
\uses{thm:cggr11-hard,lem:chain-pme}
Irreducible matrices over a field are PME exactly when they are cut-transpose equivalent,
and the list of cuts can be taken of length at most $2|I|$.
\end{theorem}
\begin{proof}
\leanok
Combine Theorem~\ref{thm:cggr11-hard} with Lemma~\ref{lem:chain-pme}.
\end{proof}

\section{The alternative Theorem 4.1}

\begin{lemma}[An active cut gives a nontrivial fiber]
\label{lem:ct-nontrivial}
\lean{Matrix.not_hasTrivialFiber_of_isCutTranspose}
\leanok
\uses{lem:ct-pme,lem:pme-basic}
If some cut-transpose of $A$ is not diagonally equivalent to $A$, the fiber of $A$ is not
trivial.
\end{lemma}
\begin{proof}
\leanok
The cut-transpose lies in the fiber by Lemma~\ref{lem:ct-pme}.
\end{proof}

\begin{theorem}[Alternative Theorem 4.1]
\label{thm:main-alt}
\lean{Matrix.hasTrivialFiber_iff_forall_isCutTranspose}
\leanok
\uses{lem:reducible-fiber,lem:ct-pme,lem:pme-basic,thm:cggr11-hard,lem:pme-irred,lem:chain-de}
For nonempty $I$ and a field $F$, the fiber of $A$ is trivial exactly when $A$ is
irreducible and every cut-transpose of $A$ at every cut of $A$ is diagonally equivalent to
$A$. No hypothesis $|I|\ge4$ is needed, since there are no cuts when $|I|\le3$.
\end{theorem}
\begin{proof}
\leanok
If the fiber is trivial, $A$ is irreducible by Lemma~\ref{lem:reducible-fiber}, and each
cut-transpose lies in the fiber by Lemma~\ref{lem:ct-pme}. Conversely, a matrix $B$ in the
fiber is irreducible by Lemma~\ref{lem:pme-irred}, so by Theorem~\ref{thm:cggr11-hard} it is
diagonally equivalent to the end of a chain starting at $A$; that end is diagonally
equivalent to $A$ by Lemma~\ref{lem:chain-de}.
\end{proof}

\begin{corollary}[No cut]
\label{cor:no-cut}
\lean{Matrix.hasTrivialFiber_of_noCut'}
\leanok
\uses{thm:main-alt}
For nonempty $I$, an irreducible matrix with no cut has a trivial fiber. Compare the
external statement Lemma~\ref{ext:no-cut}, which the proof of
Theorem~\ref{thm:cggr11-hard} itself uses.
\end{corollary}
\begin{proof}
\leanok
The cut-transpose condition is vacuous.
\end{proof}

\begin{definition}[Examples]
\label{def:ct-examples}
\lean{exMat, exActive, exActiveCT, exMatFact, exActiveFact}
\leanok
\uses{def:cut-transpose}
Over $\mathbb{Q}$, with indices $0,\dots,3$ and the cut $X=\{0,1\}$:
\[ A=\begin{pmatrix}0&2&1&1\\3&0&1&1\\1&1&0&1\\1&1&1&0\end{pmatrix},\qquad
   A^*=\begin{pmatrix}0&2&1&1\\3&0&1&1\\1&1&0&2\\1&1&3&0\end{pmatrix},\qquad
   \operatorname{ct}_X(A^*)=\begin{pmatrix}0&2&1&1\\3&0&1&1\\1&1&0&3\\1&1&2&0\end{pmatrix}. \]
Guards check that $\operatorname{ct}_X(A)=A$, that the third matrix is the cut-transpose
of the second, and that the last two have equal principal minors. A further example
checks that the second matrix of Example 4.6 (Definition~\ref{def:example46}) is the
transpose of a cut-transpose of the first.
\end{definition}
